\needspace{10\baselineskip}
\begin{center}
\pdfbookmark[1]{7. Zu Art. 90 und 100.}{zusatz-007}
[7. Zu Art. 90 und 100.]
\end{center}
\begin{center}
\begin{tabular}{c}\hline
\small Astronomisches Jahrbuch für 1814, S. 256. Berlin 1811.\\\hline
\end{tabular}
\end{center}
.... Noch füge ich Ihrem Wunsche zufolge einen kleinen Zusatz zu meiner Theoria motus corporum coelestium bei.

Zur Auflösung der wichtigen Aufgabe, aus zweien Radiis vectoribus und dem eingeschlossenen Winkel die elliptischen oder hyperbolischen Elemente zu bestimmen, habe ich mich mit grossem Vortheil einer Hülfsgrösse \(\xi\) bei der Ellipse, \(\zeta\) bei der Hyperbel bedient, für welche ich jenem Werke eine Tafel angehängt habe. Berechnet ist diese Tafel nach einem dort angeführten continuirlichen Bruche, dessen vollständige Ableitung aber dort nicht gegeben ist, und zu dessen theoretischer Entwickelung, die mit andern Untersuchungen zusammenhängt, ich bisher noch nicht Gelegenheit gefunden habe. Es wird daher manchem lieb sein, hier einen andern Weg angezeigt zu finden, auf welchem man jene Hülfsgrösse eben so bequem hätte berechnen können.

Wir haben (Art. 90)
\[
\xi=x-\tfrac56+\frac{10}{9X}=\frac{xX-\tfrac56X+\tfrac{10}{9}}{X}\mathpunct{\cdot}
\]

\SourceLeaf{307}
Der Zähler dieses Bruchs verwandelt sich leicht, wenn man für \(X\) die dort gegebene Reihe substituirt, in
\[
\begin{aligned}
&\tfrac8{105}xx\Bigl(1+\frac{2.8}{9}x+\frac{3.8.10}{9.11}xx+\frac{4.8.10.12}{9.11.13}x^3\\
&\qquad+\frac{5.8.10.12.14}{9.11.13.15}x^4+\text{etc.}\Bigr).
\end{aligned}
\]

Setzt man also die Reihe
\[
1+\frac{2.8}{9}x+\frac{3.8.10}{9.11}xx+\ldots=A,
\]
so wird
\[
\begin{aligned}
xX-\tfrac56X+\tfrac{10}{9}&=\tfrac8{105}Axx\\[.5ex]
X&=\frac{\tfrac43\bigl(1-\tfrac{12}{175}Axx\bigr)}{1-\tfrac65x}\\[.5ex]
\xi&=\frac{\tfrac2{35}Axx\bigl(1-\tfrac65x\bigr)}{1-\tfrac{12}{175}Axx}\mathpunct{\raisebox{.45ex}{,}}
\end{aligned}
\]
nach welcher Formel man \(\xi\) immer bequem und sicher berechnen kann. Für \(\zeta\) braucht man nur \(z\) anstatt \(x\) zu setzen.

Ich bemerke nur noch, dass man \(A\) noch bequemer nach folgender Formel berechnen kann
\[
\begin{aligned}
A=(1-x)^{-\frac32}\Bigl(&1+\frac{1.5}{2.9}x+\frac{1.3.5.7}{2.4.9.11}xx\\
&+\frac{1.3.5.5.7.9}{2.4.6.9.11.13}x^3+\text{etc.}\Bigr);
\end{aligned}
\]
allein die Ableitung dieser Reihe aus der vorigen beruht auf Gründen, die hier nicht ausgeführt werden können\HistoricalNoteEight{Vgl. Band III S. 209.]}.
\begin{center}\rule{.18\textwidth}{.35pt}\end{center}
